\documentclass{ip-quiz}

\quizlevel{n1}
\quiznum{3}
\quiztitle{Funciones - Opción 0}

\begin{document}

\makeheader
\maketitle

\begin{sectionbox}{Enunciado}
Defina una función que reciba como parámetro una cantidad de milisegundos totales (un entero) y retorne un \texttt{str} con el tiempo equivalente en horas, minutos, segundos y milisegundos.

\textbf{Ejemplo:} Si la función recibe \texttt{3\_723\_004}, debe retornar \texttt{"1 h, 2 min, 3 s y 4 ms"}.
\end{sectionbox}

\begin{sectionbox}{Solución}
\begin{minted}{python}
def convertir_milisegundos(milisegundos_totales: int) -> str:
    milisegundos_en_una_hora = 1000 * 60 * 60
    horas = milisegundos_totales // milisegundos_en_una_hora
    milisegundos_restantes = milisegundos_totales % milisegundos_en_una_hora

    milisegundos_en_un_minuto = 1000 * 60
    minutos = milisegundos_restantes // milisegundos_en_un_minuto
    milisegundos_restantes = milisegundos_restantes % milisegundos_en_un_minuto

    milisegundos_en_un_segundo = 1000
    segundos = milisegundos_restantes // milisegundos_en_un_segundo
    milisegundos = milisegundos_restantes % milisegundos_en_un_segundo

    return f"{horas} h, {minutos} min, {segundos} s y {milisegundos} ms"
\end{minted}
\end{sectionbox}

\end{document}
